2.1 Graphical User Interface
The graphical user interface (GUI) serves as the central workspace for the protection analysis tool. The GUI organizes the different functions into accessible sections, allowing the user to move between diagrams, component selection, protection zone analysis, and electrical calculations without needing to leave the application. The interface is designed around an interactive workspace where components can be placed, connected, and modified as the protection system is developed.

2.2 Component Library
The component library can also be expanded through the Create a component functionality button. This allows a user to create a custom component when a required component is not already included in the library. The ability to create customized components makes the tool more flexible and allows it to accommodate different protection-system configurations and project requirements.
The Component Library provides the users with a collection of electrical and protection system components that industries use that can be used to construct one-line diagrams. The Component Library allows for both commonly used and custom created components to be selected and placed into a workspace for quicker access and efficiency, instead of needing the user to manually draw each component into the schematics.

2.3 Protection Zone Outliner:
The Zone Outliner will be used to organize and visualize the different areas of Zones of Protection found within the electrical schematics. Within the sandbox, once the components and circuit connections are established, the Zone Outliner would create a display that would identify the protection zones.
The purpose of having this feature is make it easier in relationships between the schematic components and their associated protection zones to be easier to understand. Separating the electrical system into individual zones, would allow easier identification among it’s users for which components are included within a particular zone and how evaluating the different parts of the system is protected.

2.4 Momentary Calc:

2.5 PDF:
After importing the PDF, the diagram can be displayed in the background of the Sandbox while the user recreates with additional component edits/changes or analyzes the system using the available components.
The background image as well as the grid can be toggled on or off, allowing the users to switch between viewing the original diagram and viewing of the newly created components within the sandbox. 
